%% 自動生成（phys/gen_phys.py）。直接編集しないこと
\makeatletter
\phdefvol{5}{原子}{量子から半導体・原子力へ}{「LEDはなぜ光る？」「放射線はどこから出る？」に、原子で答えよう}
\phdefface{1}{飛び出す}
\phdefface{2}{すりぬける}
\phdefface{3}{切りかえる}
\phdefface{4}{分かれる}
\def\phFrontLine{スマホのチップも、X線の検査も、原子炉も。原子の4つの顔で読みとく}
\phdefpart{1}{第1部}{飛び出す：光と電子の二重性}{（原-1〜原-4）}{1}
\phdefpart{2}{第2部}{すりぬける：量子力学の入口}{（原-5〜原-7）}{2}
\phdefpart{3}{第3部}{切りかえる：固体と半導体}{（原-8〜原-10）}{3}
\phdefpart{4}{第4部}{分かれる：原子核と放射線}{（原-11〜原-16）}{4}
\def\ph@whyA{花火の色は、なぜ材料によって変わる？}\def\ph@whyrA{原-4\hspace{0.8em}\bkpageno{62}}
\def\ph@whyB{電子顕微鏡は、なぜふつうの顕微鏡より小さなものが見える？}\def\ph@whyrB{原-3\hspace{0.8em}\bkpageno{46}}
\def\ph@whyC{LEDは、どうやって電気を直接、光に変える？}\def\ph@whyrC{原-9\hspace{0.8em}\bkpageno{154}}
\def\ph@whyD{大昔の木や骨の年代は、どうやってわかる？}\def\ph@whyrD{原-12\hspace{0.8em}\bkpageno{210}}
\def\ph@whyE{X線写真で、なぜ骨が白くうつる？}\def\ph@whyrE{原-13\hspace{0.8em}\bkpageno{228}}
\def\ph@whypage{304}
\def\phBackCatch{気になる「なぜ？」から、\\ページをめくってみよう}
\makeatother
